\documentclass[../main]{subfiles}
\ifSubfilesClassLoaded{\setcounter{chapter}{7}}{}
\begin{document}

\chapter{Animation dan Animator}

\begin{subcpmk}
  \item Sub-CPMK 6.1: Membuat Animator Controller dan mengatur state serta transition
  \item Sub-CPMK 6.2: Menggunakan parameter (Trigger, Bool) dan kontrol animasi dari script
\end{subcpmk}

\section{Collision vs Trigger}

Collision dan trigger menjadi landasan mekanik game. Bab ini membahas OnTriggerEnter dan OnCollisionEnter. \class{Collider} mendeteksi benturan \cite{unityPhysics}. Dua mode:

\begin{itemize}
  \item \textbf{Collision}: Objek saling bertabrakan secara fisik (memantul, terdorong). Untuk dinding, lantai, objek padat
  \item \textbf{Trigger}: Hanya mendeteksi overlap, tidak ada respon fisik. Centang ``Is Trigger'' di Inspector. Untuk collectibles, checkpoint, zona
\end{itemize}

Untuk collectibles: gunakan \textbf{Trigger} agar bola bisa ``melewati'' dan mengumpulkan tanpa mendorong objek. Jika pakai Collision biasa, bola akan mendorong collectible.


\subsection{OnTriggerEnter vs OnCollisionEnter}

\begin{itemize}
  \item \texttt{OnTriggerEnter(Collider other)}: Dipanggil saat collider lain masuk trigger. Gunakan untuk collectibles
  \item \texttt{OnCollisionEnter(Collision collision)}: Dipanggil saat collision fisik terjadi. Gunakan untuk tabrakan dengan dinding
\end{itemize}

\subsection{OnTriggerEnter}

\begin{unitycode}
void OnTriggerEnter(Collider other)
{
    if (other.gameObject.CompareTag("PickUp"))
    {
        other.gameObject.SetActive(false);
    }
}
\end{unitycode}

\section{Animation State dan Transisi}

\begin{learningoutcomes}
\item Memahami State Machine animasi
\item Mengatur transition antar state
\item Menggunakan parameter (Trigger, Bool)
\end{learningoutcomes}

\subsection{State Machine}

State: Idle, Walk, Attack, Death. Transition: kondisi untuk berpindah state. Parameter: Trigger (one-shot), Bool (on/off), Float, Int. Contoh: parameter ``isDead'' Trigger $\rightarrow$ transition ke state Death.

\subsection{Blend Tree}

Untuk animasi smooth (contoh: blend Walk dan Run berdasarkan speed). Gunakan untuk movement blend.


\begin{aktivitas}
  \item Buat animasi explosion untuk asteroid
  \item Implementasikan animasi di script saat collision
  \item Eksplorasi Animator parameters dan transitions
\end{aktivitas}

\begin{latihan}
  \item Jelaskan perbedaan Animator Controller dan Animation Clip!
  \item Apa fungsi parameter Trigger di Animator?
  \item Bagaimana memanggil animasi dari script?
\end{latihan}

\begin{asesmen}
\textbf{Praktik}: Buat animasi explosion dan integrasikan ke Space Survivor.
\textbf{Rubrik}: Lihat Lampiran A
\end{asesmen}

\begin{checklist}
  \item Saya memahami Animator Controller
  \item Saya dapat membuat Animation Clip
  \item Saya dapat menggunakan Animator.SetTrigger
\end{checklist}

\begin{rangkuman}
Animator Controller mengatur state machine animasi. Animation Clip berisi data gerakan. Gunakan Animator.SetTrigger untuk memicu transisi. Animation Event untuk callback saat frame tertentu.
\end{rangkuman}

\ifSubfilesClassLoaded{
  \renewcommand{\bibname}{Daftar Pustaka}
  \bibliographystyle{plain}
  \bibliography{../references}
}{}
\end{document}
